\subsection{Competing branches in local search}
\label{subsec:black-box-datascience-finite-budget-optimization}

The two branch-history implementations had the two largest median
best-so-far numerical systolic ratios in this comparison. Seven fixed
implementations shared the same 64 generic ten-facet starting bodies, a cap
of 128 full evaluations and a nominal one-second accumulated
evaluator-plus-optimizer budget. The comparison uses a heuristic
floating-point objective, rather than the certified capacity implementation
of Section~\ref{sec:numerics}. Its geometry and volume are floating-point
computations.
\begin{center}
\begin{tabular}{lrr}
\hline
Implementation & Median best ratio & Across-start $10$--$90\%$ range\\
\hline
Four-anchor history & .985 & .957--.999\\
History with transition prediction & .982 & .956--.994\\
Finite-gap model & .977 & .948--.995\\
Literal gradient & .881 & .790--.924\\
Safeguarded gradient & .855 & .780--.916\\
CMA-ES & .818 & .718--.869\\
Coordinate poll & .601 & .358--.787\\
\hline
\end{tabular}
\end{center}
The lower-envelope structure of the branch model in
Section~\ref{subsec:first-order-search-model} explains why increasing one
currently minimizing branch can expose another as the new minimum. The
finite-gap method therefore models branches near the current minimum.
The history methods also retain branches from earlier points; one predicts
additional branches from transitions. A numerically chosen branch window
is not guaranteed to contain every relevant branch.

CMA-ES~\cite{dsHansen2016} and coordinate polling supply population and
direct-search comparisons. These are rankings of the fixed implementations
under this allocation, not of optimally tuned algorithm families. The ranges
in the table show variation across starting bodies; all 448 best-so-far
ratios are below one. Appendix~\ref{app:ds-local-optimization} specifies the
available implementations, traces and budget interpretation.

